\documentclass[article,11pt,oneside,a4paper,english,brazil,sumario=tradicional]{abntex2}
\usepackage[T1]{fontenc}
\usepackage[utf8]{inputenc}
\usepackage{mathptmx}
\usepackage{indentfirst}
\usepackage{graphicx}
\usepackage{microtype}
\usepackage{url}
\usepackage{enumitem}
\usepackage{booktabs}
\usepackage{array}
\usepackage{tabularx}
\usepackage[font=small,labelfont=bf]{caption}
\usepackage{needspace}
\usepackage{tikz}
\usetikzlibrary{arrows.meta,positioning}
\usepackage[abnt-emphasize=bf,abnt-and-type=e,alf]{abntex2cite}
\setsecheadstyle{\bfseries\normalsize\uppercase}
\setsubsecheadstyle{\bfseries\normalsize}
\setlrmarginsandblock{2.5cm}{2.5cm}{*}
\setulmarginsandblock{2.5cm}{2.5cm}{*}
\checkandfixthelayout
\setlength{\parindent}{1.25cm}
\setlength{\parskip}{0pt}
\SingleSpacing
\setlength{\parskip}{1.3pt}
\renewcommand{\arraystretch}{1.13}
\setlength{\tabcolsep}{7pt}
\setlength{\afterchapskip}{0pt}
\setlength{\beforechapskip}{0pt}
\setlength{\beforesecskip}{0.59\baselineskip}
\setlength{\aftersecskip}{0.27\baselineskip}
\begin{document}
\selectlanguage{brazil}
\frenchspacing
\thispagestyle{empty}
\begin{center}
\uppercase{\bfseries Influência do limiar de decisão na classificação de tumores mamários: comparação entre MLP e Regressão Logística}\\[8pt]
\uppercase{\bfseries Influence of decision thresholds on breast tumor classification: MLP versus Logistic Regression}
\end{center}
\begin{flushright}
Alane Damasceno Moreno\\
{\small Bacharelado em Ciência da Computação -- IFCE, Campus Aracati}\\
{\small Inteligência Artificial -- Prof. Gabriel Duarte}
\end{flushright}
\vspace{-0.2em}
\section{PROBLEMA DE PESQUISA}

A classificação supervisionada de tumores mamários pode apresentar elevado desempenho global e, ainda assim, deixar de identificar parte dos casos malignos. Nesses problemas, falsos negativos merecem análise específica, pois a acurácia isolada não expressa a distribuição dos erros. O conjunto Breast Cancer Wisconsin Diagnostic, disponibilizado no \textit{scikit-learn}, fornece dados públicos para experimentos controlados com 569 amostras e 30 atributos.

Srinivasu et al. (2024) investigaram uma combinação de CatBoost e MLP, enquanto Bamodu, Nezam e Chung (2026) compararam oito algoritmos e discutiram interpretabilidade e calibração de probabilidades. A pergunta deste projeto é: \textit{como a variação do limiar de decisão afeta o recall, o F1-score e os falsos negativos de uma MLP e de uma Regressão Logística, sob o mesmo protocolo experimental?}

\section{JUSTIFICATIVA}
O limiar padrão de 0,5 não necessariamente produz a melhor relação entre sensibilidade e falsos positivos em todas as tarefas. Comparar decisões obtidas a partir das mesmas probabilidades permite evidenciar a diferença entre qualidade de ordenação (ROC-AUC) e classificação em um ponto operacional. O estudo utiliza ferramentas públicas, sem pretender criar um sistema de diagnóstico clínico. A lacuna é tratada como \textit{recorte analítico entre os estudos selecionados}, não como alegação de ineditismo absoluto.

\section{OBJETIVO GERAL}
Analisar o efeito da variação do limiar de decisão sobre a identificação de tumores malignos em modelos de Regressão Logística e MLP treinados no mesmo conjunto de dados.

\section{OBJETIVOS ESPECÍFICOS}
\begin{enumerate}[leftmargin=*,itemsep=2pt,topsep=3pt]
\item Implementar pipelines equivalentes de treinamento e validação para Regressão Logística e MLP.
\item Comparar limiares de 0,1 a 0,9 quanto a falsos negativos, falsos positivos, recall e F1-score.
\item Avaliar ROC-AUC e matrizes de confusão, discutindo o compromisso entre sensibilidade e precisão.
\item Identificar limitações metodológicas e possibilidades de replicação do estudo.
\end{enumerate}

\begin{table}[!ht]
\centering
\caption{Síntese do problema e das decisões de pesquisa}
\small
\begin{tabularx}{\textwidth}{@{}>{\bfseries}p{3.4cm}X@{}}
\toprule
Dimensão & Delimitação adotada \\
\midrule
Problema & Efeito do limiar de decisão sobre a identificação de malignidade. \\
Comparação & Regressão Logística (referência) e MLP (modelo neural). \\
Foco analítico & Recall, F1-score e número de falsos negativos. \\
Escopo & Estudo acadêmico com base pública, sem aplicação clínica. \\
\bottomrule
\end{tabularx}
\par\vspace{2mm}{\footnotesize Fonte: elaboração própria (2026).}
\end{table}

\section{METODOLOGIA}
Será realizado um estudo quantitativo, aplicado e experimental, em Python no Google Colab (CPU), com \textit{scikit-learn}. Os dados públicos Breast Cancer Wisconsin Diagnostic (UCI, disponibilizados pelo \textit{scikit-learn}) serão carregados por \texttt{load\_breast\_cancer}, com 569 observações, 30 atributos e duas classes; serão verificadas as condições de licença e citação na fonte original. A codificação será ajustada para 1 = maligno (classe positiva) e 0 = benigno. A divisão treino--teste será estratificada (80/20, semente 42). O \texttt{StandardScaler} será ajustado apenas ao treino, dentro de \textit{Pipeline}, prevenindo vazamento de dados. A Regressão Logística será o \textit{baseline}; a MLP será o comparador neural. A avaliação incluirá validação cruzada estratificada de cinco partições no treinamento.

Serão obtidas probabilidades fora da amostra na validação cruzada para descrever os limiares de 0,1 a 0,9. O critério exploratório de seleção será o maior F1 da classe maligna no treino, com desempate pelo menor limiar; o conjunto de teste permanecerá reservado para uma única avaliação final dessa regra. Serão apresentadas matriz de confusão, falsos negativos/positivos, recall, precisão, F1-score e ROC-AUC. Será comparado o limiar padrão ao selecionado, sem interpretação clínica.

\begin{table}[!ht]
\centering
\caption{Protocolo experimental previsto}
\small
\begin{tabularx}{\textwidth}{@{}>{\bfseries}p{3.4cm}X@{}}
\toprule
Componente & Procedimento \\
\midrule
Base e divisão & Wisconsin Diagnostic: 569 amostras, 30 atributos; treino/teste estratificado 80/20 (semente 42). \\
Pré-processamento & Padronização dentro do \textit{Pipeline}, ajustada apenas aos dados de treinamento. \\
Modelos & \texttt{LogisticRegression} e \texttt{MLPClassifier}. \\
Validação & Cinco partições estratificadas nos dados de treino. \\
Análise & Limiares de 0,1 a 0,9; matriz de confusão, precisão, recall, F1 e ROC-AUC. \\
Ambiente & Python, \textit{scikit-learn} e Google Colab (CPU). \\
\bottomrule
\end{tabularx}
\par\vspace{2mm}{\footnotesize Fonte: elaboração própria (2026).}
\end{table}


\begin{figure}[!ht]
\centering
\begin{tikzpicture}[>=Stealth, font=\footnotesize,
 stage/.style={draw=black!55, rounded corners=2pt, minimum height=0.55cm, align=center, inner sep=4pt}]
\node[stage, text width=2.5cm] (a) {Base Wisconsin\\569 amostras};
\node[stage, text width=2.5cm, right=0.28cm of a] (b) {Treino/teste\\estratificado};
\node[stage, text width=2.5cm, right=0.28cm of b] (c) {Pipeline\\LR e MLP};
\node[stage, text width=2.5cm, right=0.28cm of c] (d) {Limiares e\\métricas};
\draw[->, thick] (a)--(b);
\draw[->, thick] (b)--(c);
\draw[->, thick] (c)--(d);
\end{tikzpicture}
\caption{Fluxo previsto para o experimento comparativo.}
\par\vspace{1mm}{\footnotesize Fonte: elaboração própria (2026).}
\end{figure}

\section{RESULTADOS ESPERADOS}
Espera-se produzir um notebook reproduzível, tabelas comparativas e gráficos que mostrem como o limiar altera os erros e as métricas dos dois classificadores. O êxito será definido pela execução reproduzível do protocolo e por uma discussão fundamentada das diferenças observadas, independentemente de qual algoritmo obter a maior pontuação.

\section{RISCOS E LIMITAÇÕES}
A amostra é relativamente pequena, utiliza um único conjunto público e não representa diversidade de contextos clínicos. A seleção de limiar por validação interna pode variar com as partições; portanto, os resultados não serão generalizados para prática médica. Se a MLP apresentar instabilidade, será adotada configuração previamente definida.

\section{DECLARAÇÃO DE USO DE INTELIGÊNCIA ARTIFICIAL}
O Hermes Agent, conectado à OpenRouter, foi utilizado para auxiliar na localização e triagem inicial de referências científicas. Ferramentas de IA também auxiliaram na revisão linguística, incluindo correção ortográfica, coesão gramatical e ajustes estilísticos, e o apoio à organização lógica dos argumentos, com atenção à integridade acadêmica e sob supervisão e revisão da autora. As sugestões incorporadas ao texto foram examinadas e revisadas pela autora, que assume integralmente a responsabilidade pelo conteúdo final, pelas escolhas metodológicas, pelos argumentos e pelas conclusões apresentados. Não foram compartilhadas informações inéditas, dados sensíveis ou conteúdos originais de terceiros com as plataformas utilizadas. Esses cuidados foram adotados visando ao respeito aos direitos autorais e à proteção de dados pessoais, considerando a Lei Geral de Proteção de Dados Pessoais (LGPD --- Lei nº 13.709/2018). Declaro que li e conferi pessoalmente as referências citadas neste documento, verificando sua identificação bibliográfica e a correspondência entre as afirmações e as fontes originais. 

\clearpage
\section*{REFERÊNCIAS}
\begingroup\footnotesize\setlength{\parskip}{0pt}
\noindent SRINIVASU, P. N. et al. XAI-driven CatBoost multi-layer perceptron neural network for analyzing breast cancer. \textit{Scientific Reports}, v. 14, art. 28674, 2024. DOI: 10.1038/s41598-024-79620-8.\\
BAMODU, O. A.; NEZAM, S.; CHUNG, C.-C. Explainable machine learning for breast cancer prediction in resource-constrained settings. \textit{PLOS Digital Health}, v. 5, n. 9, e0001706, 2026. DOI: 10.1371/journal.pdig.0001706.\\
PEDREGOSA, F. et al. Scikit-learn: Machine Learning in Python. \textit{JMLR}, v. 12, p. 2825--2830, 2011.\\
SCIKIT-LEARN DEVELOPERS. \textit{Breast cancer Wisconsin (diagnostic) dataset}. Disponível em: \url{https://scikit-learn.org/stable/modules/generated/sklearn.datasets.load_breast_cancer.html}. Acesso em: 22 out. 2026.

\endgroup
\end{document}